\documentclass[11pt]{article}
\usepackage[margin=0.85in,top=0.75in,bottom=0.85in]{geometry}
\usepackage{amsmath}
\usepackage{amssymb}
\usepackage{enumitem}
\usepackage{hyperref}
\usepackage{titlesec}

\setlist{itemsep=1pt, topsep=3pt, parsep=0pt}
\setlength{\parskip}{2pt}
\titlespacing*{\section}{0pt}{10pt}{4pt}
\titlespacing*{\subsection}{0pt}{7pt}{3pt}

\title{\textbf{BloodLine -- Predictive Blood Shortage \& Smart\\ Donor Matching Platform}\\[0.3em]
\large Project Proposal for UCS503P\\
Submitted to: Dr. Anushka\\
Thapar Institute of Engineering and Technology}

\author{
Avraj Singh Sekhon, CSED\thanks{Roll No: 1024030287, Email: asekhon\_be24@thapar.edu} \and
Krish Kumar, CSED\thanks{Roll No: 1024030296, Email: kkumar3\_be24@thapar.edu} \and
Kyna Sood, CSED\thanks{Roll No: 1024030303, Email: ksood\_be24@thapar.edu} \and
Manya Sabherwal, CSED\thanks{Roll No: 1024030288, Email: msabherwal\_be24@thapar.edu}
}

\date{August 18, 2026}

\begin{document}
\maketitle

\section{Higher-order goal \dots secondary framing}
This project aims to improve blood availability by reducing the time and coordination effort required to identify and respond to potential blood shortages. While the topic is broader than any single blood bank, the immediate engineering objective is to translate proactive blood availability management into a measurable, deployable software solution.

\section{Time-to-value \dots primary heuristic}
\begin{itemize}
\item We will deliver meaningful increments quickly by prototyping the core user workflow of inventory tracking, demand forecasting, shortage-risk detection, and donor matching within a short iteration window.
\item We will prioritize fast feedback over perfection by using weekly iterations (and targeting CI-backed automation early) to reduce release friction.
\item Success is defined by what can be delivered and measured in the first iteration, then improved based on testing and user feedback.
\end{itemize}

\section{Problem Statement}
Blood banks and hospitals can face uncertainty and delays in maintaining adequate blood availability. Common pain points include:
\begin{itemize}
\item \textbf{Reactive shortage management}: current inventory can be monitored, but short-term future shortage risk requires forecasting and early warning.
\item \textbf{Fragmented coordination}: inventory, nearby blood-bank availability, hospital requirements, and potential donors are not unified into one workflow.
\item \textbf{Inefficient donor outreach}: finding suitable nearby donors based on blood group, donation history, proximity, availability, and consent can be time-consuming.
\end{itemize}

\noindent\textbf{Why this matters now (contextual relevance):} A proactive, data-driven coordination layer can help blood-bank staff prepare for potential shortages earlier and reduce avoidable delays in donor outreach or cross-bank coordination.

\section{Proposed Solution \dots Software Proposition}
\subsection{Overview}
Build a web-based ``BloodLine'' system with the following components:
\begin{itemize}
\item \textbf{Inventory dashboard}: authorized blood-bank staff record and monitor blood stock by blood group/component.
\item \textbf{Demand forecasting}: estimate near-term demand from historical data, initially using a simulated dataset where suitable granular real data is unavailable.
\item \textbf{Shortage-risk detection}: compare projected demand with available/expected inventory and flag potential shortages.
\item \textbf{Smart donor matching}: identify potential nearby donors using configured blood-group, cooldown, proximity, availability, and consent rules.
\item \textbf{Role-based workflow}: provide separate interfaces for donors, blood-bank managers, inventory staff, medical/screening staff, hospital coordinators, and system administrators.
\end{itemize}

\subsection{Core workflow}
\begin{enumerate}
\item Blood-bank staff update current inventory and donation records.
\item BloodLine forecasts near-term demand and flags potential shortage risk.
\item The manager reviews the risk, nearby availability, and potential donor matches.
\item Donor outreach is initiated; medical staff perform final screening and the resulting donation is recorded.
\item Recording a donation triggers an inventory update, which re-runs the shortage-risk check so the manager's dashboard reflects whether the flagged risk has been resolved or still requires further outreach.
\end{enumerate}

\subsection{Operational constraints for an educational, multi-user setting}
\begin{itemize}
\item Keep the solution usable on low-to-mid range devices via responsive web design.
\item Use clearly labelled simulated historical demand data where granular real consumption data is unavailable for the prototype.
\item BloodLine does not make final medical eligibility decisions; donor screening and actual redistribution remain with authorized blood-centre staff.
\item Nearby blood-bank availability and cross-bank data sharing are simulated/assumed for the MVP, since real integration would require inter-organizational data-sharing agreements outside the scope of this course project; the data layer is designed to accept authorized real feeds later.
\end{itemize}

\section{Solution Approach \dots Engineering focus}
This is an engineering course project, so we focus on reliable multi-user workflow, role-based access, system correctness, and deployability rather than novel algorithm research.

\subsection{Forecasting and donor-matching assistance \dots non-research}
Demand forecasting will use a time-series approach such as Prophet or SARIMA. The model will estimate near-term demand from the available historical/simulated dataset. A transparent rule-based layer will compare projected demand with available inventory and generate a shortage risk signal. Potential donors will then be filtered using blood-group compatibility, donation cooldown, approximate location, availability, and notification consent. Final medical eligibility remains with qualified staff.

The simulated demand dataset will be constructed to reflect realistic consumption patterns rather than arbitrary random values: a weekly seasonal cycle (e.g., lower elective-surgery demand on weekends), periodic demand spikes (e.g., festival/holiday periods and simulated emergency events), and injected random noise per blood group/component. This keeps the forecasting task non-trivial and prevents the model from simply learning the exact generating function.

\subsection{Web architecture \dots web-based solution}
A typical 3-tier web architecture:
\begin{itemize}
\item \textbf{Frontend}: responsive UI for donors, blood-bank staff, hospitals, and administrators.
\item \textbf{Backend API}: authentication, role-based authorization, inventory operations, forecasting, risk detection, matching, requests, and notifications.
\item \textbf{Database}: store users, roles, blood banks, inventory, donations, demand history, predictions, requests, and alerts.
\end{itemize}

\subsection{CI/CD and fast delivery enabler}
CI/CD will be used to:
\begin{itemize}
\item Automatically build and run tests on every change.
\item Produce repeatable deployment artifacts to staging.
\item Enable safe and frequent releases of incremental features.
\end{itemize}

\section{Evaluation Criterion \dots measurable and attributable}
We define success using measurable metrics tied to the core workflow.

\subsection{Primary evaluation metric}
\textbf{Forecasting error}: compare predicted demand with actual values in the held-out test period using MAE and RMSE.
\begin{itemize}
\item \textbf{Attribution}: calculated from forecast outputs and corresponding observed demand values in the test dataset.
\item \textbf{Baseline comparison}: Prophet/SARIMA results will be compared against a simple moving-average/naive forecast baseline on the same test period, so that MAE/RMSE improvements are demonstrably attributable to the modelling approach rather than to the characteristics of the simulated data.
\end{itemize}

\subsection{Secondary evaluation metrics}
\begin{itemize}
\item \textbf{Shortage-risk performance}: measure whether known shortage-risk cases in the test data are correctly identified.
\item \textbf{Donor matching correctness}: verify that configured blood-group, cooldown, proximity, availability, and consent filters are applied correctly.
\item \textbf{Workflow response time}: measure response time for common operations such as inventory update, dashboard loading, and donor search.
\item \textbf{Notification targeting}: verify that alerts are limited to donors satisfying the configured matching rules.
\end{itemize}

\subsection{Pilot validation plan \dots high-level}
\begin{itemize}
\item Use a simulated demand dataset and divide it chronologically into training and testing periods.
\item Compare metrics before/after within the iteration window using forecast outputs and seeded donor/blood-bank scenarios.
\end{itemize}

\section{Scalability \dots with foresight}
The proposition should scale in both theoretical and practical senses.
\begin{enumerate}
\item \textbf{Scalable (theoretically)}: The system should support growth in number of blood banks, donors, inventory records, and concurrent dashboard usage by:
\begin{itemize}
\item decoupling components (frontend/backend),
\item enabling pagination and indexing for common queries,
\item using stateless API design.
\end{itemize}
\item \textbf{Operationally deployable (practically)}: The system should run on common web hosting:
\begin{itemize}
\item managed database,
\item platform-hosted deployment,
\item CI/CD pipeline for staging/production.
\end{itemize}
\end{enumerate}

\section{Engine availability heuristic}
A mature set of ``engines'' (tooling/services) is available to implement this as a web-based project:
\begin{itemize}
\item React.js + Tailwind CSS for responsive web interfaces.
\item FastAPI (Python) for REST APIs and application logic.
\item PostgreSQL for structured relational data.
\item Pandas with statsmodels/Prophet for demand forecasting.
\item Google Maps API or equivalent for proximity and regional visualization.
\item Authentication, notifications, testing framework, CI runner, and staging deployment target.
\end{itemize}
We will use these mature components to focus on workflow design, correctness, integration, and measurement rather than inventing new infrastructure.

\section{Project scope and deliverables \dots iteration-friendly}
\subsection{Initial deliverable \dots first iteration}
\begin{itemize}
\item Role-based login and dashboards.
\item Blood-bank inventory and donor records.
\item Simulated demand dataset and forecasting pipeline.
\item Shortage-risk calculation and basic donor matching.
\item Basic logging and timestamps.
\item CI: build + backend unit tests + linting; CD to staging.
\end{itemize}

\subsection{Subsequent deliverables}
\begin{itemize}
\item Regional blood-availability heatmap and nearby surplus/shortage coordination.
\item Targeted in-app/push notification workflow.
\item Hospital blood-request and medical screening workflows.
\item Improved forecasting evaluation and indexed queries for scale readiness.
\end{itemize}

\section{Risks and mitigations}
\begin{itemize}
\item \textbf{Data availability}: use clearly labelled simulated data for the MVP and keep the data layer ready for authorized real data.
\item \textbf{Forecast uncertainty}: present predictions as risk signals rather than guarantees and retain human review.
\item \textbf{Medical/privacy constraints}: final donor clearance remains with qualified staff; minimize personal data and restrict access by role.
\item \textbf{Operational issues}: start with in-app/push notifications and use staging early for repeatable deployments.
\end{itemize}

\section{Summary}
This project proposes a deployable web platform to make blood availability more proactive by combining inventory tracking, short-term demand forecasting, shortage-risk detection, nearby availability, smart donor matching, and role-based coordination. It meets the course criteria by emphasizing time-to-value through rapid iterations, CI/CD, measurable evaluation, and scalability readiness. It remains focused on engineering workflow and system correctness rather than research-driven novelty.

\end{document}
